\begin{corollary}\label{sga3:XXI:3.1.10}
If \(\Delta\) is a system of simple roots, then
\(\operatorname{ind}(\Delta^\ast)\) is a system of simple coroots
\emph{(i.e.}\ a system of simple roots of \(\cR^\ast\)\emph{)}.
\end{corollary}

\begin{proof}
If
\[
 \beta=\sum_{\alpha\in\Delta}a_\alpha\alpha,
 \qquad a_\alpha\in\mathbb N,
\]
then
\[
 p(\beta)=\sum_{\alpha\in\Delta}a_\alpha p(\alpha).
\]
Consequently, by \eqref{sga3:XXI:eq:10},
\[
 \beta^\ast
 =\sum_{\alpha\in\Delta}
   a_\alpha\frac{\ell(\alpha)}{\ell(\beta)}\alpha^\ast,
\]
which shows that \(\operatorname{ind}(\Delta^\ast)\) satisfies
condition (i\('\)) of Remark~\ref{sga3:XXI:3.1.9}.
\end{proof}

By Proposition~\ref{sga3:XXI:2.1.1}, if \(\alpha^\ast\) is not
indivisible, then
\[
 \alpha^\ast=2u^\ast,
 \qquad
 u^\ast\in\operatorname{ind}(\Delta^\ast)
 \quad\text{and}\quad
 u=2\alpha.
\]
It follows that:\footnote{Editors' note: The preceding sentence was
added, and in Corollary 3.1.11 the expression
\(4a_\alpha\ell(\alpha)\) was replaced by
\(2a_\alpha\ell(\alpha)\).}

\begin{corollary}\label{sga3:XXI:3.1.11}
Let \(\Delta\) be a system of simple roots, and let
\[
 \beta=\sum_{\alpha\in\Delta}a_\alpha\alpha,
 \qquad a_\alpha\in\mathbb Z,
\]
be the expression of \(\beta\) relative to \(\Delta\). Then
\(\ell(\beta)\) divides \(2a_\alpha\ell(\alpha)\); it even divides
\(a_\alpha\ell(\alpha)\) if \(\alpha^\ast\) is indivisible.
\end{corollary}

Before continuing to state the properties of systems of simple roots,
we show that such systems exist.

\subsection{Systems of positive roots}

\begin{definition}\label{sga3:XXI:3.2.1}
A subset \(R_+\subset R\) is called a \emph{system of positive roots}
of \(\cR\) \emph{(or of \(R\); see
Remark~\ref{sga3:XXI:3.2.2}\emph{)}} if it satisfies:
\begin{enumerate}
\item[(P 1)] \(R_+\cap(-R_+)=\varnothing\);
\item[(P 2)] \(R_+\cup(-R_+)=R\);
\item[(P 3)] \((\Qp\cdot R_+)\cap R\subset R_+\).
\end{enumerate}
In particular, such a set is closed. We shall see later, in
Corollary~\ref{sga3:XXI:3.3.8}, that a closed set satisfying
(P 1) and (P 2) also satisfies (P 3), and hence is a system of
positive roots. If \(w\in W\) and \(R_+\) is a system of positive
roots, then \(w(R_+)\) is a system of positive roots.
\end{definition}

\begin{remark}\label{sga3:XXI:3.2.2}
This definition involves only \(R\). We shall also say that \(R_+\)
is a system of positive roots of \(R\).
\end{remark}

\begin{remark}\label{sga3:XXI:3.2.3}
Conditions (P 1) and (P 2) immediately give
\[
 \operatorname{Card}(R_+)=\frac{\operatorname{Card}(R)}{2}.
\]
It follows that if \(R_+\) and \(R'_+\) are two systems of positive
roots and \(R_+\subset R'_+\), then \(R_+=R'_+\).
\end{remark}

\begin{remark}\label{sga3:XXI:3.2.4}
If \(R_+\) is a system of positive roots, then \(R^\ast_+\) is a
system of positive coroots \emph{(i.e.}\ a system of positive roots
of \(R^\ast\)\emph{)}.
\end{remark}

This follows immediately from Lemma~\ref{sga3:XXI:1.1.4} and
Scholium~\ref{sga3:XXI:1.2.2}.

\begin{definition}\label{sga3:XXI:3.2.5}
Let \(\Delta\) be a system of simple roots. Set
\[
 \mathscr{P}(\Delta)
 =(\Qp\cdot\Delta)\cap R
 =(\mathbb N_+\cdot\Delta)\cap R.
\]
\end{definition}
